\pagebreak
\section{Метод решения}
Токенизатор написан на C++17 со стандартной библиотекой (STL для токенизации разрешена правилами курса). Функция \texttt{tok::tokenize} вынесена в заголовочный файл \texttt{src/tokenizer.h}, чтобы подключить её к индексатору в следующих работах; \texttt{src/main.cpp} --- утилита командной строки \texttt{build/tokenizer}. Сборка --- \texttt{make}.

Входные данные --- корпус из ЛР1: 1\,000\,000 текстов документов (\texttt{lab1/data\_1m/text/<источник>/*.txt}: 620\,164 статьи Wikipedia, 296\,250 постов Reddit, 48\,768 вопросов Travel Stack Exchange с ответами, 34\,818 статей Wikivoyage; английский язык, UTF-8, 5,6~ГиБ).

\subsection{Правила токенизации}
\begin{enumerate}
    \item Текст декодируется из UTF-8 в символы Юникода. Некорректная последовательность байтов считается разделителем.
    \item Токен --- максимальная непрерывная последовательность букв и цифр. Буква или цифра --- символ, для которого \texttt{iswalnum} истинно в локали \texttt{C.UTF-8} (классы символов glibc по базе Юникода): латиница, в т.\,ч. с диакритикой, кириллица, греческий, арабский, иероглифы, цифры.
    \item Все остальные символы --- разделители: пробельные (в т.\,ч. неразрывный пробел), знаки препинания, дефисы и тире, апострофы, кавычки, скобки, \texttt{\& + / @ \% °}, комбинируемые диакритические знаки, мягкий перенос.
    \item Регистр понижается: в ASCII --- сдвигом кода, у остальных символов --- функцией \texttt{towlower} (\texttt{É} $\to$ \texttt{é}, \texttt{Ж} $\to$ \texttt{ж}, \texttt{İ} $\to$ \texttt{i}). Других нормализаций нет.
    \item Длина токена считается в символах Юникода, а не в байтах.
\end{enumerate}
Пример: \texttt{Lviv's Old Town (Ukrainian: Старе Місто) --- São Paulo, 3.14, B-52} $\to$ \texttt{lviv s old town ukrainian старе місто são paulo 3 14 b 52}.

Текст обрабатывается за один проход: очередной символ декодируется и, если это буква или цифра, в нижнем регистре дописывается в текущий токен; на разделителе накопленный токен передаётся функции обратного вызова. Символы ASCII проверяются и переводятся в нижний регистр сравнениями, без вызова функций библиотеки.

\subsection{Утилита командной строки}
\begin{itemize}
    \item \texttt{tokenizer < text > tokens} --- текст со стандартного ввода, токены по одному на строку в стандартный вывод. Используется в автотестах и для анализа токенов корпуса.
    \item \texttt{tokenizer -s < files > stats} --- на вход список путей к файлам, по одному на строку. Каждый файл читается целиком и токенизируется. Выводятся число файлов, байтов, токенов, средняя длина токена, общее время, отдельно время чтения файлов и время токенизации, скорость в мкс/КиБ и КиБ/с. Время измеряется внутри программы (\texttt{std::chrono::steady\_clock}).
\end{itemize}

Замеры выполняет скрипт \texttt{scripts/bench.py}: статистика по источникам и время на первых $N$ документах корпуса ($N$ от 1000 до 1\,000\,000, документы источников чередуются, пока источник не кончится). Перед замерами корпус читается один раз (статистика по источникам), чтобы файлы попали в кэш ОС. Каждый замер повторяется 3 раза, берётся минимальное время.
